%!TEX root = ../QuantumNoiseAndDecoherence.tex
% ──────────────────────────────────────────────────────────────
%  Lecture 3 — Stinespring Dilation and Fock Space
% ──────────────────────────────────────────────────────────────

\section{Complete positivity and dilation theory}%
\label{sec:cp-dilation}

In the previous lecture we defined quantum channels axiomatically:
a channel $\chan\colon \algB \to \alg$ must be unital and completely
positive.  We now ask: \emph{can every channel be built from simple,
physically transparent operations?}

% ──────────────────────────────────────────────────────────────
\subsection{Constructive approach to quantum operations}%
\label{sec:constructive-ops}

We enumerate three elementary operations that should certainly be
allowed in quantum mechanics.

\subsubsection{Adjoining an ancilla}

Given a system~$A$ in state~$\omega$, we may bring in an additional
system~$B$ in a fixed state~$\sigma$:
\begin{equation}\label{eq:adjoin}
  \omega \;\longmapsto\; \omega \tp \sigma\,.
\end{equation}
In the Schr\"odinger picture, this takes a state of~$A$ to a state
of~$AB$.

\subsubsection{Unitary evolution}

A closed system undergoes unitary dynamics:
\begin{equation}\label{eq:unitary-op}
  \rho \;\longmapsto\; U\,\rho\, U^\adj\,,\qquad U^\adj U = \One\,.
\end{equation}

\subsubsection{Reduction to a subsystem}

If we have a composite system~$AB$ and decide to ignore~$B$
(never measure anything on~$B$), the effective state of~$A$ is
obtained by partial trace:
\begin{equation}\label{eq:reduce-op}
  \rho_{AB} \;\longmapsto\; \tr_B(\rho_{AB})\,.
\end{equation}

\begin{intuition}[Building arbitrary operations]
  By composing these three elementary operations, we can build
  remarkably complex processes.  For example: start with a
  qubit~$A$; adjoin an ancilla~$B$ in state~$\sigma$; apply a
  joint unitary~$U$ on~$AB$; adjoin a third system~$C$ in
  state~$\sigma_C$; apply a different unitary~$V$ on~$BC$;
  finally discard~$A$, leaving an effective state of~$BC$.  Each
  step is one of the three allowed operations.
\end{intuition}

% ──────────────────────────────────────────────────────────────
\subsection{The Stinespring dilation theorem}%
\label{sec:stinespring}

The remarkable fact is that, for finite-dimensional systems, the
axiomatic and constructive approaches are \emph{equivalent}.

\begin{theorem}[Stinespring dilation]\label{thm:stinespring}
  Let $\chan\colon \Mn{N} \to \Mn{M}$ be a completely positive, unital
  map (i.e.\ a quantum channel in the Heisenberg picture).  Then there
  exists a positive integer~$L$ and an isometry
  $V\colon \C^M \to \C^N \tp \C^L$ (satisfying $V^\adj V = \One_M$)
  such that
  \begin{equation}\label{eq:stinespring}
    \eqbox{\chan(X) = V^\adj\,(X \tp \One_L)\,V}
    \qquad \text{for all } X \in \Mn{N}\,.
  \end{equation}
  The integer~$L$ can be chosen so that $L \leq M^2 N^2$.
\end{theorem}

In the Schr\"odinger picture, the dual map $\chan_*$ acts on states as:
\begin{equation}\label{eq:stinespring-schrodinger}
  \chan_*(\rho) = \tr_L\!\bigl(V\,\rho\, V^\adj\bigr)\,.
\end{equation}
This is precisely a composition of the three allowed operations:

\begin{center}
\begin{tikzpicture}[
  block/.style={rectangle, rounded corners=3pt, draw=cgblue,
    fill=cgblue!8, minimum height=2em, minimum width=2.2cm,
    font=\sf\footnotesize, text=spacecadet, inner sep=5pt},
  arr/.style={-{Stealth[length=5pt]}, thick, cgblue}
]
  \node[block] (s1) {$\rho$};
  \node[block, right=0.8cm of s1] (s2) {$\rho \tp \ketbra{0}{0}$};
  \node[block, right=0.8cm of s2] (s3) {$U(\rho \tp \ketbra{0}{0})U^\adj$};
  \node[block, right=0.8cm of s3] (s4) {$\tr_L(\cdots)$};
  \draw[arr] (s1) -- (s2)
    node[midway, above, font=\tiny]{adjoin};
  \draw[arr] (s2) -- (s3)
    node[midway, above, font=\tiny]{unitary};
  \draw[arr] (s3) -- (s4)
    node[midway, above, font=\tiny]{reduce};
\end{tikzpicture}
\end{center}

\begin{remark}
  The Stinespring theorem is only valid for \emph{finite-dimensional}
  systems.  For infinite-dimensional systems, there exist completely
  positive, unital maps that cannot be written in the
  form~\eqref{eq:stinespring}.  This means there may be legal quantum
  evolutions that do not arise from ``system~$+$~environment~$+$~unitary''---a
  fact with implications for, e.g., black hole information loss.
\end{remark}

% ══════════════════════════════════════════════════════════════
\section{Quantum many-particle systems}\label{sec:many-particle}

We now turn to modeling the \emph{environment} in our
system--environment picture.  Real environments---the electromagnetic
field, phonons in a solid, a thermal bath---consist of many (possibly
indefinitely many) particles.  We need mathematical machinery to handle
this.

% ──────────────────────────────────────────────────────────────
\subsection{Fock space}\label{sec:fock-space}

\subsubsection{Motivation}

We already know how to compose~$n$ quantum systems with Hilbert
spaces $\hilb_{A_1},\ldots,\hilb_{A_n}$: the joint Hilbert space is
\begin{equation}\label{eq:n-fold-tensor}
  \hilb_{A_1} \tp \hilb_{A_2} \tp \cdots \tp \hilb_{A_n}\,.
\end{equation}
But many physical systems---a mode of the electromagnetic field, Cooper
pairs in a superconductor---have an \emph{indefinite} number of
particles: particles can be created and destroyed.  For such systems,
we need a Hilbert space that accommodates \emph{any} number of
particles simultaneously.

\subsubsection{Definition}

Let $\hilb$ be the single-particle Hilbert space.  The \textbf{Fock
space} (for distinguishable particles) is
\begin{equation}\label{eq:fock-space}
  \eqbox{\mathcal{F}(\hilb)
    = \bigoplus_{n=0}^{\infty} \hilb^{\tp n}}\,,
\end{equation}
where $\hilb^{\tp 0} = \C$ (the ``vacuum sector'') and
$\hilb^{\tp n} = \underbrace{\hilb \tp \cdots \tp \hilb}_{n}$.

A vector $\Psi \in \mathcal{F}(\hilb)$ is a sequence
\begin{equation}\label{eq:fock-vector}
  \Psi = (\alpha_0,\; \ket{\alpha_1},\; \ket{\alpha_2},\; \ldots)\,,
  \qquad
  \alpha_0 \in \C,\;\;
  \ket{\alpha_n} \in \hilb^{\tp n}\,,
\end{equation}
with the requirement that $\sum_n \norm{\alpha_n}^2 < \infty$.

\begin{remark}
  One might ask: why a direct sum rather than an infinite tensor
  product?  The infinite tensor product $\hilb^{\tp\infty}$ is strictly
  larger and contains states with no physical interpretation (e.g.\
  superpositions involving infinitely many particles in incompatible
  ways).  Fock space is the \emph{smallest} Hilbert space that assigns
  a vector to every state we expect to exist.
\end{remark}

The distinguished vector
$\ket{\Omega} = (1, 0, 0, \ldots) \in \mathcal{F}(\hilb)$
is the \textbf{vacuum state}: zero particles.

% ──────────────────────────────────────────────────────────────
\subsection{Identical particles and symmetrisation}%
\label{sec:identical-particles}

For \emph{identical} particles (photons, electrons, \ldots), swapping
two particles must leave the state unchanged (up to a sign).  We
enforce this by projecting onto the appropriate symmetry sector.

\subsubsection{The symmetric group and permutation operators}

The \textbf{symmetric group}~$S_n$ is the group of all permutations of
$\{1,2,\ldots,n\}$.  Each permutation $\pi \in S_n$ acts on product
vectors by reordering the factors:
\begin{equation}\label{eq:perm-action}
  U_\pi\,(\ket{v_1} \tp \cdots \tp \ket{v_n})
  = \ket{v_{\pi^{-1}(1)}} \tp \cdots \tp \ket{v_{\pi^{-1}(n)}}\,.
\end{equation}

Define the \textbf{symmetrisation} and \textbf{antisymmetrisation}
projectors:
\begin{equation}\label{eq:sym-projectors}
  P_n^{\pm}
  = \frac{1}{n!}\sum_{\pi \in S_n}
    \varepsilon_\pm(\pi)\, U_\pi\,,
\end{equation}
where $\varepsilon_+(\pi) = 1$ for all~$\pi$ (bosons) and
$\varepsilon_-(\pi) = \mathrm{sgn}(\pi)$ is the sign of the
permutation (fermions).

\subsubsection{Bose and Fermi Fock space}

\begin{definition}[Bose/Fermi Fock space]\label{def:bose-fermi-fock}
  The \textbf{Bose Fock space} and \textbf{Fermi Fock space} are
  \begin{equation}\label{eq:bose-fermi-fock}
    \eqbox{\mathcal{F}_\pm(\hilb)
      = \bigoplus_{n=0}^{\infty} P_n^{\pm}\,\hilb^{\tp n}}\,.
  \end{equation}
\end{definition}

In this course, we focus almost exclusively on \emph{bosons}
($\mathcal{F}_+$), since our environments will typically be
electromagnetic fields (photons).  Everything we develop, however,
carries over to fermions with only sign changes.

\begin{example}[Single-mode Fock space]\label{ex:single-mode}
  If $\hilb = \C$ (a single mode of the electromagnetic field, e.g.\
  the lowest mode of an optical cavity), then
  $\hilb^{\tp n} \cong \C$ for every~$n$, and Fock space becomes
  \begin{equation}\label{eq:single-mode-fock}
    \mathcal{F}_+(\C) = \bigoplus_{n=0}^{\infty} \C
    \cong \ell^2(\N)\,.
  \end{equation}
  This is the Hilbert space of a single quantum harmonic oscillator,
  spanned by number states $\ket{0}, \ket{1}, \ket{2}, \ldots$ The
  space is infinite-dimensional even though the single-particle space
  is one-dimensional---there can be any number of photons in the mode.
\end{example}

\begin{example}[Multi-mode cavity]\label{ex:multi-mode}
  If a cavity supports $k$ modes (e.g.\ two polarisations of the
  lowest frequency), then $\hilb = \C^k$ and Fock space
  $\mathcal{F}_+(\C^k)$ describes an indefinite number of photons,
  each of which can occupy any of the $k$~modes.
\end{example}

% ──────────────────────────────────────────────────────────────
\subsection{Creation and annihilation operators}%
\label{sec:creation-annihilation}

Given a single-particle state $\ket{f} \in \hilb$, we define operators
on the (distinguishable-particle) Fock space $\mathcal{F}(\hilb)$
that add or remove a particle in state~$\ket{f}$.

\begin{definition}[Annihilation and creation operators]%
\label{def:annihilation-creation}
  The \textbf{annihilation operator} $a(f)$ and \textbf{creation
  operator} $a^\adj(f)$ act on the $n$-particle sector as:
  \begin{align}
    a(f)\,\ket{\alpha_n}
    &= \sqrt{n}\;(\bra{f} \tp \One^{\tp(n-1)})\,\ket{\alpha_n}
    \in \hilb^{\tp(n-1)}\,,
    \label{eq:annihilate}
    \\[4pt]
    a^\adj(f)\,\ket{\alpha_n}
    &= \sqrt{n+1}\; P_{n+1}^\pm\,(\ket{f} \tp \ket{\alpha_n})
    \in \hilb^{\tp(n+1)}\,,
    \label{eq:create}
  \end{align}
  and $a(f)\,\ket{\Omega} = 0$ (one cannot annihilate from the vacuum).
\end{definition}

When restricted to the Bose or Fermi Fock space via the
symmetrisation projectors~$P_n^\pm$, these operators satisfy the
fundamental \textbf{canonical commutation/anticommutation relations}:
\begin{equation}\label{eq:ccr-car}
  \eqbox{[a(f),\, a^\adj(g)]_\mp = \braket{f}{g}\,\One}\,,
  \qquad
  [a(f),\, a(g)]_\mp = 0\,,
\end{equation}
where $[\cdot,\cdot]_-$ is the commutator (bosons) and
$[\cdot,\cdot]_+$ is the anticommutator (fermions).

\begin{keyresult}[The number operator]
  For a single mode ($\hilb = \C$, $\ket{f} = 1$), writing
  $a \coloneqq a(1)$, the \textbf{number operator}
  $\hat{n} = a^\adj a$ satisfies
  $\hat{n}\ket{n} = n\ket{n}$, where
  $\ket{n} = \frac{(a^\adj)^n}{\sqrt{n!}}\ket{0}$
  are the familiar number states of the quantum harmonic oscillator.
\end{keyresult}

\begin{remark}
  The notation $a(f)$ emphasises that the annihilation operator depends
  \emph{linearly} on the single-particle state~$f$, while $a^\adj(f)$
  depends \emph{antilinearly} on~$f$.  In much of the physics
  literature, one fixes a mode basis $\{e_k\}$ and writes
  $a_k \coloneqq a(e_k)$, recovering the familiar mode-indexed ladder
  operators.
\end{remark}
